% ============================================================
% sections/discussion.tex  考察（初稿 2026-10-26）
% 方針：成功も否定的な結果も「シミュレーションで他者を理解する仕組みの，できることとできないこと」として置く。
%   推測は推測と書く。発達研究は類推にとどめる。版の番号は使わない。
% ============================================================
\section{考察}\label{sec:discussion}

\subsection{共有された予測の仕組みで見られたもの，見られなかったもの}\label{sec:disc-summary}
本研究は，自他で重みを共有する予測の仕組みが，自分の価値の経験から他者について何を生み，何を生まないかを調べた．
結果は次のようにまとめられる．

見られたものは三つある．
第一に，与えた快・不快の信号のもとで価値に従って動いた経験から，報酬も目標の入力も用いない予測学習によって，価値に従う行動の見込みが生まれ，その強さは経験の量に応じて段階的に強まった（実験1）．
ただし，この経験の効果は学習データの位置の分布の変化と分けられていない．
第二に，その見込みは，他者のための内部状態に入れた快・不快にも働いた（実験2）．
自他の重みの共有だけを外すと，自分側の見込みはそのままで他者側の働きだけが消えた．
このモデルでは，自分側で学んだ重みが他者側の計算に使われることが，見込みを他者に運ぶために要る．
ただし，この対照の結果は構成からほぼ予想されるものであり，重ね合わせ機構という仕組みに固有の効果であることまでは示していない．
第三に，運ばれた見込みは限られた形で使える．
研究者が用意した二つの好みの候補を入れ比べて予測の当たりを比べると，経験80\%条件では偶然より高く好みの方を選べた（実験4）．
ただし主の自己視界基準での正解率は，他者の進んだ方角だけで選ぶ基準を上回らず，副の他者特徴基準でも組をまたいでわずかに上回るにとどまった（組ごとにはっきり上回ったのは1組）．

調べた条件で確認できなかったものも三つある．
第一に，他者用モジュールは，構成を4通りに変えても他者の好みを書き出さなかった（実験3）．
第二に，見込みは他者に合わせて調整されずに当てはめられ，行動の仕方が自分と違う他者では，本当の快・不快を与えたときの予測が，経験の少ない条件より外れた（実験2の真値条件．段階IIを行った4組のうち，緑好みデータで2組，好み混合データで3組では，何も与えない場合より悪かった）．
第三に，照合による見分けのうち，他者が進んだ方角では説明できない部分は一部の組に限られた（実験4）．

つまりこの仕組みは，他者を「価値の方へ動く者」として予測する枠組みを，自分の経験から作って他者に当てはめる．
しかし，調べた条件では，その枠組みの中で，この他者が何を好み，どれだけ強く好みに従うかを読み取る働きは確認できなかった．
シミュレーションによる他者理解の，できることとできないことが，一つのモデルの中で分かれて現れたと言える．

\ref{app:switch}の好み切替条件（A-1 の好きな色がエピソードごとに変わる経験）では，自分側の見込みは好きな色ごとに対称に生まれたが，他者側への当てはめと照合は強くならなかった．
見込みを他者に運ぶには，経験の量だけでなく，一貫した好みに従う経験であることが効いている可能性がある．

\subsection{投影の誤りとシミュレーション説}\label{sec:disc-projection}
経験80\%条件で真値条件の予測が悪くなったことは，シミュレーション説が予想する自己中心的な誤り\cite{goldman2006simulating}の，計算モデルの上での一つの現れと読める．
この解釈では，モデルは，快が高ければほぼ好きな色の角にいた自分の経験から，他者も同じ強さで快の方へ動くと見積もり，半分しか好みに従わない A-2 を本当より快の方へ寄せて予測した（\ref{sec:res-projection}節）．
追加の比較では，A-2 も自分と同じ割合で好みに従うとき，経験80\%条件で本当の快・不快がより役立つようになり，その変化は経験の少ない条件より大きかった（\ref{sec:res-projection}節）．
これは，悪化が自他の行動の仕方の食い違いに依るという考えを支持し，学習で見ていない入力であることだけによる説明とは合わない（ただし組によっては悪化が残った）．
ただし，この比較では価値に従って動いた経験と角に多くいた経験とを分けられないので，以下の議論は投影という解釈が正しい場合のものである．
ここで投影されているのは「何を好むか」ではなく，「どれだけ強く好みに従うか」という自分の行動の仕方である．
A-1 自身の好きな色への寄りは，調べた経験50\%条件と好み切替条件のモデルでは見られなかった（経験80\%条件では調べていない）．

人の他者理解では，自分を出発点にして，そこから他者に合わせて調整する過程が想定されている\cite{epley2004perspective,tamir2010neural}．
また Suzuki ら\cite{suzuki2012learning}は，自分の価値づけを流用する信号と，模擬から予測した他者の行動と実際の行動との差を表す信号とを区別し，後者による補正は他者の価値が自分と異なるときに重要になると示唆した．
本研究の段階Iのモデルは，この区別でいえば，自分の価値づけを流用する経路だけを持ち，他者の行動の予測の外れから他者の価値を補正する経路を持たない．
真値条件の悪化は，調整を欠いたシミュレーションが，自分と違う他者にどう失敗するかを示す例になりうる．
ただし，これらの研究が述べるのは人の判断や脳の信号についてであり，本モデルの構成要素がそれらと一対一に対応するとは主張しない．

他者用モジュールは，他者側の予測の外れを減らすように学習するので，Suzuki らの行動の予測の外れに当たる学習信号は与えられていた．
実際，経験80\%条件の上で学習したモジュールの出力を入れると，他者側の特徴予測の誤差は 0.112〜0.188 で，本当の快・不快を入れた場合（0.376〜0.418）よりずっと小さかった（組1〜4，二つのデータ，表\ref{tab:projection}）．
補正そのものは学ばれたのである．
しかし，その補正は他者の価値を直す形ではなく，予測に役立つ位置や方角の手がかりを伝える形をとった（\ref{sec:res-module}節）．
% 出典：paper/figs/projection_values.json（numbers_check.md #20）。v13 まとめ §4 の 0.12〜0.15・0.40〜0.42 は組1〜3・緑好み・ep200 の値。

好み切替条件では，当てはめが弱くなると同時に投影も見られなかった．
これは，投影が一貫した自分の好みから生まれるという解釈と合う．
ただし，当てはめそのものが弱まったために投影も見えなくなった可能性もあり，二つは区別できていない．
% 出典：v14 まとめ §3，執筆メモ §5 付録実験A。

\subsection{他者用モジュールはなぜ好みを書き出さないか}\label{sec:disc-module}
文献\cite{noguchi2022superposition}では，同じ構成の Motion Generator が process-2 の運動の枠を埋めるよう学習し，その出力は，A-1 が A-2 と同じように動いたときの運動に似た値になった．
自他の対応に乗って，他者の観察から「自分ならそう動いたときの運動」が作られたのである．
本研究の他者用モジュールは，同じ手順で快・不快の枠を埋めるよう学習したが，好みを書き出さなかった．
この違いは，次のように解釈できる（推測）．

文献\cite{noguchi2022superposition}の課題では，他者の運動は他者の次の位置をほぼ決めるので，枠に本当の運動を書くことが，そのまま予測に最も役立つ．
自分に対応する信号を書くことと，予測をよくすることが一致している．
本研究の課題では，A-2 は一歩ごとに半分しか好みに従わないので，本当の快・不快は A-2 の動きを部分的にしか決めない．
実際，本当の快・不快を入れたときの予測の改善は，経験50\%条件でも小さく，経験80\%条件では誤差が経験の少ない条件より大きくなった．
一方，位置や方角の手がかりを伝える出力は，予測をずっとよくした．
ここで「近道」とは，学習の目的（予測の外れを減らすこと）はよく満たすが，意図した表現（他者の価値）とは別の手がかりに頼る解を指す．
深層学習では，学習データでの目的はよく満たすが，条件が変わると通用しない規則を学んでしまう現象が shortcut learning として論じられており\cite{geirhos2020shortcut}，ここでの近道もそれに相当する可能性がある．ただし，本研究では，条件を変えたときにこの近道が通用しなくなることまでは確かめていない．
モジュールは予測の外れだけで学習するので，好みを書く道ではなく，位置や方角を書く近道が選ばれたと考えられる．

モジュールの作りを変えた三つの試みも，この解釈と合う．
出力の形を限った形制限モジュールは，範囲 2〜5 を通り道にすることを防いだが，範囲ごとの強さで位置を伝えた．
他者遮蔽学習は視界から A-2 の位置を取りにくくしたが，好みは書かれなかった．
1ビットモジュールは書ける量を1つの符号に絞ったが，符号は A-2 の進む方角に沿った．
抜け道をふさいでも好みが書かれなかったことについては，好みを伝えても予測の改善が小さいことが一つの説明として考えられる．
ただし，これを直接操作して確かめてはいない．
この説明が正しければ，投影の誤りとモジュールの近道は同じ原因から生まれることになる．
自分の行動の仕方と他者の行動の仕方が違うので，自分の見込みに他者の本当の快・不快を入れても予測が当たらず，モジュールにはそれを書く理由がない．

一方で，実験4のように，候補ごとの予測を観察と比べる手続きを外から与えると，同じ見込みから好みをある程度選べた（図\ref{fig:write-choose}）．
見込みの中には好みを区別する手がかりがあるが，予測の外れを減らすだけの学習では，それを取り出す方へ進まなかったと言える．
より一般的に言えば，予測の誤差だけを目的にすると，同じくらい予測をよくする表し方が複数あるとき，意図した表し方（他者の価値）が選ばれるとは限らない．
本研究の結果は，その一つの例と位置づけられる．
ただし，内部状態については十分に調べていない．経験50\%条件の事後の分析では，$h^2$ から A-2 の好みを読み出せたが，経験の有無で差がなかった．他者用モジュールの内部状態は調べていない．
出力に好みが表れないことは，内部に好みの情報がないことを意味しない．
また，結論は調べた構成，損失，課題に限られ，別の損失や学習の条件で好みが出力されるかは分からない．

\subsection{発達研究との対応}\label{sec:disc-dev}
本研究の結果は，発達研究の知見と次のように対応づけられる．
ただし，これは計算モデルの振る舞いと発達の現象との類推であり，人間について何かを示すものではない．

Sommerville ら\cite{sommerville2005action}は，自分で物をつかむ経験を先にした3か月児が，その後に見た他者の行為の目標に注意を向けることを示した．
本研究で，価値に従って動いた経験の量が自分側の見込みを強め，それが共有を通して他者側の働きと照合の成否を変えたことは，自分の行為の経験が他者の行為の理解を変えるという方向で，この知見と対応する．

Repacholi と Gopnik\cite{repacholi1997early}の14か月児は，実験者が自分とは違う食べ物を好むことを示しても，自分の好む食べ物を差し出した．
本研究で，見込みが他者に合わせて調整されずに当てはめられ，他者用モジュールが他者固有の好みを書き出さなかったことは，この自己中心的な反応と形が似ている．
ただし，乳児が当てはめたのは自分の好みの中身であるのに対し，本モデルが当てはめたのは好みに従う強さであり，対応は部分的である．
また，本研究の結果は，こうした調整がどのような経験や仕組みから生まれるかを示すものではない．

\subsection{限界}\label{sec:limitations}
本研究には次の限界がある．

第一に，人間のデータはなく，発達研究との対応は類推にとどまる．
第二に，環境は小さい．
単一のアリーナに四隅のランドマークを置き，カメラは回転しない．
そのため「赤の方へ動く」ことは「赤の象限の方へ動く」ことと区別しにくい．
また，A-2 は一歩ごとに半分の確率でランダムに動くので，他者用モジュールや照合の結論は，この課題の作りに依存する．
第三に，視覚エンコーダは基準エンコーダに固定した．
自他を同じ符号で表す性質は学んだものではなく与えたものであり，自他の対応と色の表現を両立して学ぶエンコーダは得られていない（\ref{app:encoder}）．
第四に，快・不快は人が方角ごとに計算して与えた値であり，範囲 2〜5 は常に0である．
価値そのものをネットワークが作ったわけではなく，本研究が示したのは，与えた価値の信号とそれに従う行動の規則性を，予測学習が利用できるようになることまでである．
第五に，実験4の候補は研究者が用意した二択であり，照合の手続きも外から与えた．
ネットワークが自ら他者の好みを推し量ったとは言えない．
第六に，経験の量の操作は，学習データの位置の分布も大きく変えた．
経験80\%条件の学習データでは A-1 が赤の角にとどまり，赤から距離5以内にいるフレームの割合は約0.85であった．
本研究の経験の効果（実験1・2・4）と真値条件の悪化（実験2）は，価値に従って動いた経験の効果と，特定の場所に多くいた経験の効果とを分けられていない．
第七に，経験80\%条件のモデルは，事前登録した前提である自己位置の追跡の基準を6組とも満たさなかった．
実験1・2・4 の指標はモデル自身の予測から測るので計算はできるが，経験80\%条件の値は，共通テストの場面に十分に一般化していないモデルのものである．
第八に，非共有条件の対照では，他者側の快・不快を受ける重みが学習されないので，他者側の働きが消えることは構成からほぼ予想される．
学習中の共有以外の方法（学習後に重みを写すなど）で同じ働きが得られるかは調べておらず，効果が重ね合わせ機構に固有かは分からない．
第九に，独立に学習した組は本番の3組，追加を含めても6組と少なく，組ごとのばらつきも大きい．
最終5点は互いに相関するチェックポイントであり，独立した繰り返しではない．
指標は多いが多重比較の補正は行っていない．
また，経験の量は3水準だけで，強まり方の形は分からない．
第十に，実験4では，経験なし条件のモデルでも好みを区別できる組があった．
組5では経験なし条件の BA 自体が 0.5 を上回り，組3では経験なし条件も食い違いの条件を満たした．
条件を満たすことだけでは経験の効果とは言えず，主の判定は経験なし条件との差で行った．
第十一に，他者側の段階的な強まりの根拠の一部，投影の誤りの解釈，他者用モジュールが好みを書かない理由の分析は，結果を見た後の指標と推測に基づく．

\subsection{今後の課題}\label{sec:future}
最初の課題は，本研究の主張の原因を確かめる対照を加えることである．
一つ目に，位置の分布をそろえて，価値と行動の結びつきだけを変える対照が要る．
例えば，経験80\%条件と同じ軌跡のデータで，快・不快の符号だけをエピソードごとに無作為にすれば，位置の分布はまったく同じまま，「快の方へ動く」という結びつきだけがなくなる．
この対照で自分側・他者側の見込みや真値条件の悪化が消えれば，同じ軌跡と位置の分布のもとでも，価値の信号と行動の対応がそれらに寄与していたと言える．ただし，これは価値の信号と行動の対応に限った操作であり，赤へ近づくという行動の履歴そのものは変わらない．符号を無作為にした信号の大きさや時間変化も位置の情報を含みうるので，快・不快を常に零にした対照も並べるとよい．
二つ目に，\ref{sec:disc-module}節の解釈からは，他者の行動の仕方が自分と同じなら，他者用モジュールは価値を書き出すようになるはずだという予測が出る．
自分と同じ強さで好みに従う他者では，自分の見込みに本当の快・不快を入れることが予測に役立ち，モジュールにはそれを書く理由が生まれるからである．
このうち，本当の快・不快が役立つようになるという前半は，追加の比較で支持された（ただし組によっては悪化が残った．\ref{sec:res-projection}節）．
モジュールが価値を書き出すかという後半を確かめるには，A-2 が 0.8 の割合で好みに従うデータで段階IIの学習を行う必要がある．
三つ目に，好みの順位でしか他者の行動を予測できない課題では，位置や方角の近道が使えなくなる．
例えば，好きな物に着くとそれが価値を失い（灰色に変わり），次に好きな物へ向かう環境では，一つ目の物を得た直後にどこへ向かうかは，それまでの進む方角からは決まらず，好みの順位でしか決まらない．
このような場面で，照合が方角では説明できない見分けを示すか，また他者用モジュールが価値を書き出すかを調べる．
ただしこの環境では，観察から順位を推す材料が本当にあるか（それまでの行動から次の向かう先をどれだけ当てられるか）を，学習の前にデータだけで確かめる必要がある．

そのほか，次の課題がある．
非共有条件のモデルで，学習の後に自分側の重みを他者側へ写して他者側の働きが戻るかを調べ，効いているのが学習中の共有か，自分側で学んだ重みの再利用かを切り分ける．
他者用モジュールと共有モジュールの内部状態から，位置と方角をそろえた場面で好みが読み出せるかを調べ，出力に表れないことと内部に情報がないことを区別する．
経験の量を3水準より細かく変え，自分側の見込み，他者側の働き，照合の正解率がどう変わるかを調べる．
他者の行動の予測の外れから他者の価値を補正する経路を加え，投影の誤りが減るかを調べる．
自他の対応と価値の表現を同時に学ぶエンコーダを作り，自他の対応を与えずに同じ結果が得られるかを確かめる．
候補を二択で与える代わりに，Raileanu ら\cite{raileanu2018modeling}のように他者の価値の入力を連続的に最適化する手続きと組み合わせ，照合を研究者の用意した候補に頼らない形に広げる．
また，Sommerville らの研究との対応をより直接に調べるため，自分で動いた経験と，同じ軌跡を受け身で経験した場合とを比べることも考えられる．
